\section{Finite transfer and conductor descent on singular links}
\label{sec:supplied-p126}

Let $(R,\mathfrak m)$ be an excellent reduced Noetherian local ring,
let $\mathfrak a$ be a proper $\mathfrak m$-primary ideal, and let
$R\to B$ be finite with $B\ne0$ reduced. This section follows the supplied
manuscript \cite{SuppliedP12620260924}. Its operators act on continuous
valuative functions and explicit ideal-generated cones. They do not
identify those cones with an independently geometric singular psh class.

Use the primary sections $V_{\mathfrak a}(R)$ of
\cref{sec:local-excellent-boundary-links}; on $B$ use the extended ideal
$\mathfrak aB$. Semivaluations with nonzero prime support are retained.
For a finitely generated ideal $I$, write $g_I(v)=v(I)$. Ideal arithmetic
 gives $g_{IJ}=g_I+g_J$ and $g_{I+J}=\min(g_I,g_J)$, including infinite
values. A gauge of an ideal containing a power of the boundary ideal is
finite and continuous on its section.

\subsection{Restriction and fiber extrema}

\begin{lemma}\label{lem:supplied-p126-finite-restriction}
Restriction is a continuous closed surjection
\[
 f:V_{\mathfrak aB}(B)\longrightarrow
 Z:=\{v\in V_{\mathfrak a}(R):v(\ker(R\to B))=\infty\}.
\]
In particular it is a quotient map onto its image.
\end{lemma}

\begin{proof}
Compactness is the closed-product argument of
\Cref{thm:p121-excellent-primary-links}, which applies equally to the
finite $R$-algebra $B$ and the finitely generated ideal $\mathfrak aB$.
Restriction is continuous and its image lies in $Z$. To lift $v\in Z$,
put $\mathfrak p=\operatorname{supp}(v)$. Lying over for the integral
injection $R/\ker(R\to B)\hookrightarrow B$ supplies a prime
$\mathfrak P$ over $\mathfrak p$. The extension
$\operatorname{Frac}(B/\mathfrak P)/\operatorname{Frac}(R/\mathfrak p)$
is finite. Extend the valuation ring of $v$ to a dominating valuation
ring in this finite field extension, as in
\Cref{lem:p121-normalization-semivaluation-lift}.
An algebraic equation shows that a positive integer multiple of each
new value lies in the original real value group: among its terms at
least two attain the minimal value. Dividing by that integer gives a
well-defined order-preserving real-valued extension. Integral elements
have nonnegative value, since a negative value would make the leading
term in a monic equation uniquely smallest. Restore $\mathfrak P$ with
value $\infty$ to obtain a lift on $B$. Its boundary gauge is one.
Compactness and the Hausdorff property now give the closed quotient map.
\end{proof}

For $H\in C(V_{\mathfrak aB}(B),\mathbf R)$ define, on the exact image $Z$,
\[
 f_{\min}H(v)=\min_{w\in f^{-1}(v)}H(w),\qquad
 f_{\max}H(v)=\max_{w\in f^{-1}(v)}H(w).
\]
The first function is lower semicontinuous and the second upper
semicontinuous. Indeed,
$\{f_{\min}H\leq t\}=f(\{H\leq t\})$ is closed, and the analogous
superlevel identity applies to the maximum. Both operations preserve
order, contract the uniform norm, commute with nonnegative scaling and
addition of a pulled-back base function, and recover that base function
from its pullback. They compose in finite towers: minimize or maximize
first on each intermediate fiber and then on the intermediate points
in the image of the upper map. No finite value is assigned outside the
image of the composite.

\subsection{Field-norm coefficient ideals}

For the next construction let $A$ be an excellent normal domain and
$A\subset B$ a finite domain extension. Put $K=\operatorname{Frac}(A)$,
$L=\operatorname{Frac}(B)$ and $n=[L:K]$. If $J$ is a $B$-ideal, choose
finite $A$-module generators $b_1,\ldots,b_r$ and put
\[
 P_J(T)=N_{L(T)/K(T)}\Bigl(\sum_iT_ib_i\Bigr),\qquad
 N_{B/A}(J)=(\text{coefficients of }P_J)\subset A.
\]
The norm in this definition is a field norm; it does not require the
$A$-module $B$ to be free.

\begin{proposition}\label{prop:supplied-p126-norm-coefficients}
The coefficients belong to $A$, and $N_{B/A}(J)$ is independent of the
chosen generators. Let $\Sigma$ be the set of $K$-embeddings of $L$ in
an algebraic closure, put $q=n/|\Sigma|$, and choose a finite field
extension $M/K$ containing all their images. Let $C$ be the integral
closure of $A$ in $M$. For every nonnegative real semivaluation $v$ of
$A$, choose an extension $w$ to $C$ and set $w_\sigma=w\circ\sigma$ on
$B$. Then
\[
 g_{N_{B/A}(J)}(v)=q\sum_{\sigma\in\Sigma}g_J(w_\sigma).
 \tag{70.7f}\label{eq:supplied-p126-norm-gauge}
\]
This includes semivaluations with nonzero support and infinite ideal
values. In particular, as identities of gauges,
\[
 g_{N_{B/A}(IJ)}=g_{N_{B/A}(I)}+g_{N_{B/A}(J)},\qquad
 g_{N_{B/A}(IB)}=n g_I
\]
for $B$-ideals in the first formula and an $A$-ideal in the second.
\end{proposition}

\begin{proof}
Over an algebraic closure $\Omega$ of $K$, the algebra
$L\otimes_K\Omega$ is a product of Artinian local algebras, indexed by
$\Sigma$, each of length $q$ and residue field $\Omega$. Multiplication
by an element on a factor is its residue scalar plus a nilpotent map.
Consequently
\[
 P_J(T)=\prod_{\sigma\in\Sigma}
                  \Bigl(\sum_iT_i\sigma(b_i)\Bigr)^q.
\]
Each coefficient is integral over $A$ and belongs to $K$; normality puts
it in $A$. Replacing generators by $A$-linear combinations substitutes
linear forms in $P_J$, so the new coefficient ideal is contained in the
old one. Substitution in both directions proves independence.

Excellence makes $C$ finite, and
\Cref{lem:supplied-p126-finite-restriction} supplies $w$; its proof does
not require a local base before imposing a primary section. The minimum
valuation of the coefficients of a polynomial is multiplicative. To
verify this with prime support, first quotient by the support and pass
to the valued fraction field. Divide each nonzero polynomial by a
coefficient of least value. Its reduction over the residue field is
nonzero, and the product of two such reductions is nonzero. This proves
multiplicativity. If a polynomial vanishes modulo the support, the
corresponding values are infinite and the same identity holds.
Apply this identity to the product above. The coefficient minimum of
its $\sigma$-factor is $\min_iw(\sigma(b_i))=g_J(w_\sigma)$, proving
\eqref{eq:supplied-p126-norm-gauge}. The last two formulas follow from
ideal-gauge arithmetic and restriction to $A$.
\end{proof}

The multiplicities $q$ are generic inseparable lengths. Replacing them
by ramification index times residue degree at an arbitrary real
valuation can lose valuation defect and is not part of this formula.

\subsection{Continuous positive trace over a normal target}

Assume now that $(A,\mathfrak m)$ is normal local and excellent, and
normalize by a proper primary ideal $\mathfrak a$. Write
$X=V_{\mathfrak a}(A)$ and $Y=V_{\mathfrak aB}(B)$.

\begin{lemma}\label{lem:supplied-p126-gauge-density}
The rational linear span of $g_J$, for $B$-ideals $J\supseteq\mathfrak a^N B$
with some $N$, is uniformly dense in $C(Y,\mathbf R)$.
\end{lemma}

\begin{proof}
These finite continuous gauges include the constant
$g_{\mathfrak aB}=1$ and separate points: for $b\in B$ the gauge
$g_{(b)+\mathfrak a^N B}(w)=\min(w(b),N)$ distinguishes different values
when $N$ is sufficiently large. After a common denominator, every
rational linear combination has the form $(g_I-g_J)/d$ by products,
powers and the zero gauge $g_B$. Moreover,
\[
 \min(g_I-g_J,g_K-g_H)=g_{IH+KJ}-g_{JH}.
\]
All ideals on the right contain a boundary power. The real linear span
is a separating vector lattice containing constants; the lattice
Stone--Weierstrass argument makes it dense. Explicitly, match values
at pairs of points with affine combinations of separating gauges,
then take finite minima and maxima using compactness. Approximate the
finitely many real coefficients by rationals to obtain the same closure
for the rational span.
\end{proof}

\begin{theorem}\label{thm:supplied-p126-normal-trace}
With the notation of \Cref{prop:supplied-p126-norm-coefficients},
\[
 \mu_v=q\sum_{\sigma\in\Sigma}\delta_{w_\sigma},\qquad
 T_{B/A}H(v)=\int_YH\,\rd\mu_v
 \tag{70.7g}\label{eq:supplied-p126-trace}
\]
is independent of all choices. The measure has mass $n$ and support
exactly the restriction fiber over $v$. The operator
$T_{B/A}:C(Y,\mathbf R)\to C(X,\mathbf R)$ is positive and bounded,
with norm $n$ when $X\ne\varnothing$. It satisfies
\[
 T_{B/A}1=n,\qquad T_{B/A}(f^*h)=nh,\qquad
 T_{B/A}g_J=g_{N_{B/A}(J)}
\]
for every ideal $J$ containing a boundary power. No flatness or
separability hypothesis is imposed.
\end{theorem}

\begin{proof}
By \eqref{eq:supplied-p126-norm-gauge}, the integral of every generating
gauge against $\mu_v$ is independent of the choices. All candidate
measures have mass $n$, so \Cref{lem:supplied-p126-gauge-density} gives
equality on every continuous test function and hence equality of the
finite-support measures.
Each summand lies over $v$. Conversely, given a point $u$ of that fiber,
choose the common field $M$ with the identity embedding of $B$, and
extend $u$ to its integral closure $C$. This is possible because $C$ is
integral and finite over $B$. The identity summand is $u$ and has
positive weight. Choice independence proves that every point of the
fiber has positive weight; in particular there are at most $n$ points.

For $J\supseteq\mathfrak a^N B$, adjoining any $a_0\in\mathfrak a^N$
as a redundant generator shows that $N_{B/A}(J)$ contains $a_0^n$.
Finite generation then implies that it contains a power of
$\mathfrak a$. Its gauge is finite and continuous. Formula
\eqref{eq:supplied-p126-norm-gauge}, density and the common mass bound
therefore prove continuity of $T_{B/A}H$ for every continuous $H$.
Positivity and the displayed identities follow from the fiber measures.
The norm bound is attained on $1$ whenever $X$ is nonempty.
\end{proof}

Let $\mathcal P_{\mathfrak a}(A)$ be the cone of finite nonnegative real
combinations of ideal gauges containing a boundary power; define
$\mathcal P_{\mathfrak aB}(B)$ in the same way. Trace preserves these
cones, and pullback preserves them by extension of ideals. This is an
explicit ideal-generated class.

For a finite tower $A\subset B\subset D$ of excellent normal domains,
with $A$ local, the traces compose:
\[
 T_{D/A}=T_{B/A}\circ T_{D/B}.
 \tag{70.7h}\label{eq:supplied-p126-tower}
\]
To prove this, choose one finite extension of $\operatorname{Frac}(A)$
containing every embedding of $\operatorname{Frac}(D)$, and extend the
base semivaluation once. Group the top embeddings by their restriction
to the intermediate field. Every intermediate embedding extends, and
both field degrees and separable embedding counts multiply. Thus the
inseparable multiplicities multiply as well, and the resulting Dirac
sum is the iterated fiber measure. Choice independence proves the
formula. The proof also applies to a normal semilocal intermediate
ring on the common pulled-back section, using its finitely many local
center pieces. In this domain setting the degree-normalized averages
also compose, since the degrees multiply.

\begin{example}\label{ex:supplied-p126-nonflat-normal-trace}
For $\operatorname{char}k\ne2$, put
\[
 A=k[x^2,xy,y^2]_{(x^2,xy,y^2)},\qquad B=k[x,y]_{(x,y)}.
\]
The sign-invariant ring $A$ is normal: an element of its fraction field
integral over $A$ is in the normal ring $B$, and invariance puts it in
$A$. It is excellent, and $B/A$ is finite of generic degree two.
The residue space $B/\mathfrak m_AB$ has basis $1,x,y$, so the finite
module is not flat: a finite flat module over a local ring is free,
whereas its minimal generator number is three and its generic rank two.
For $J=(x,y)B$, use the $A$-module generators $x,y,x^2,xy,y^2$.
Writing $E=T_3x^2+T_4xy+T_5y^2$, the norm polynomial is
$E^2-(T_1x+T_2y)^2$, whose coefficient ideal is $\mathfrak m_A$.
Thus $T_{B/A}g_J=g_{\mathfrak m_A}$ on the whole section, including
semivaluations with nonzero support.
\end{example}

\subsection{The exact conductor-compatible domain}

Return to the finite map $R\to B$ at the start of the section. Write
the finite normalizations as $S_R=\prod_iA_i$ and $S_B=\prod_jD_j$,
where the factors are normal domains, and let $q_R:W_R\to X$ and
$q_B:W_B\to Y$ be their primary-link quotients. A source component $D_j$
is active over $A_i$ if its minimal prime contracts to the minimal prime
defining $A_i$. The induced inclusion $A_i\subset D_j$ is finite.
Components with positive-codimension image are called vertical here;
no generic degree is assigned to them. On normalization data set
\[
 (T_f^\nu H)|_{A_i}=\sum_{j\text{ active over }i}T_{D_j/A_i}(H|_{D_j}),
 \tag{70.7i}\label{eq:supplied-p126-branch-trace}
\]
with an empty sum equal to zero. The common-section proofs above apply
to these normal semilocal factors. This defines a positive bounded
operator. Fiber extrema still detect vertical components.

Define $\mathcal D_f\subset C(Y,\mathbf R)$ to consist of those $h$ for
which $T_f^\nu(q_B^*h)$ is constant on every fiber of $q_R$.
For such an $h$ define $\operatorname{Tr}_fh$ by
$q_R^*(\operatorname{Tr}_fh)=T_f^\nu(q_B^*h)$.

\begin{theorem}\label{thm:supplied-p126-conductor-trace}
The space $\mathcal D_f$ is a closed vector subspace, and
$\operatorname{Tr}_fh$ is the unique continuous descent of the branch
trace. Its compatibility equalities occur exactly on the normalization
fibers described by the conductor quotient in
\Cref{thm:p121-excellent-conductor-descent}.
For a finite tower $R\to B\to D$, if $h\in\mathcal D_{D/B}$ and
$\operatorname{Tr}_{D/B}h\in\mathcal D_{B/R}$, then
$h\in\mathcal D_{D/R}$ and
\[
 \operatorname{Tr}_{D/R}h
   =\operatorname{Tr}_{B/R}(\operatorname{Tr}_{D/B}h).
\]
Direct admissibility alone is not an intermediate descent hypothesis.
\end{theorem}

\begin{proof}
Fiber constancy is a closed linear condition for the uniform norm.
The compact Hausdorff quotient $q_R$ makes it necessary and sufficient
for a unique continuous descent. A finite-valued conductor point has
one normalization lift; on the infinite-conductor stratum the exact
fibers are those of restriction through $S_R/\mathfrak c_R$ to
$R/\mathfrak c_R$, by \Cref{thm:p121-excellent-conductor-descent}.

A minimal component of $D$ dominating a minimal component of $R$ passes
through a unique minimal component of $B$ dominating that component:
incomparability in an integral extension makes a prime over a minimal
target prime minimal in the source. Conversely such an active chain
is active for the composite. Summing \eqref{eq:supplied-p126-tower}
over these chains gives the identity on normalizations. The two assumed
descents and uniqueness yield the displayed identity downstairs and
prove direct admissibility.
\end{proof}

Define $\mathcal P_{\mathfrak a}^{\nu}(R)$ to be the continuous functions
whose pullbacks to all normal components lie in their ideal-generated
cones. This includes both the branchwise ideal condition and the fiber
equalities. Equation \eqref{eq:supplied-p126-branch-trace} gives
\[
 \operatorname{Tr}_f\bigl(\mathcal P_{\mathfrak aB}^{\nu}(B)
                         \cap\mathcal D_f\bigr)
                  \subseteq\mathcal P_{\mathfrak a}^{\nu}(R).
\]
This normalization-ideal cone is not identified with the singular nef
or analytic psh cone.

There is a useful sufficient criterion for realization by an actual
ideal. If $S$ is the normalization, $\mathfrak c$ its conductor,
$J\supseteq\mathfrak c$ an $S$-ideal and $K\subset R/\mathfrak c$ an
ideal with $J/\mathfrak c=K(S/\mathfrak c)$, let $I$ be the inverse
image of $K$ in $R$. Then $IS=J$: both contain $\mathfrak c$ and have
the same image modulo it. Conversely every $R$-ideal containing the
conductor has this description after extension. Fiber-compatible
continuous gauges need not be gauges of individual downstairs ideals.
For the cusp $k[t^2,t^3]_{(t^2,t^3)}\subset k[t]_{(t)}$, the gauge of
$(t)$ on the maximal-ideal section is $1/2$, whereas every proper
nonzero downstairs ideal has value at least one. The former is still
the real multiple $g_{\mathfrak m}/2$; this example distinguishes a
single ideal from the real cone.

\subsection{The obstruction and the finite-flat endpoint}

\begin{example}\label{ex:supplied-p126-conductor-obstruction}
Let
\[
 \begin{gathered}
 R=k[x,y,z]_{(x,y,z)}/(xy),\qquad
 S=k[x,z]_{(x,z)}\times k[y,z]_{(y,z)},\\
 J=(x,z)\times(y,z^2).
 \end{gathered}
\]
Normalize by $\mathfrak m=(x,y,z)$. On the first branch take monomial
semivaluations $u_t(x)=t$, $u_t(z)=1$; on the second take
$w_t(y)=t$, $w_t(z)=1$, for integers $t\geq2$. Their restrictions both
converge pointwise to the semivaluation supported on $(x,y)$ and inducing
$\operatorname{ord}_z$ on $R/(x,y)$. Indeed, transverse terms acquire
unbounded value and local denominators have value zero. But
$g_J(u_t)=1$ and $g_J(w_t)=2$. On either generic branch the degree-one
trace must equal this source value. No continuous function on the target
can have both limits at the common conductor point. Thus a finite
birational surface map need not admit a continuous degree trace on all
source ideal gauges.
\end{example}

\begin{proposition}\label{prop:supplied-p126-finite-flat-trace}
If $R\to B$ is finite locally free of constant rank $n$, then
$\mathcal D_f=C(Y,\mathbf R)$. Its trace is the continuous rank-$n$
determinant trace. For a $B$-ideal $J$ containing a boundary power, its
image is the gauge of the coefficient ideal of
\[
 \det\Bigl(m_{\sum_iT_ib_i}:B[T]\longrightarrow B[T]\Bigr),
\]
where the $b_i$ generate $J$ as an $R$-module.
\end{proposition}

\begin{proof}
The determinant exists and is basis independent. Linear substitution
in both directions proves generator independence of its coefficient
ideal. Flat going down shows that every minimal source component
dominates a minimal target component. After localization at a minimal
prime of the reduced base, the reduced finite algebra splits into the
fraction fields of those components. Its determinant polynomial is
the product of their field norms. The coefficient-minimum argument in
\Cref{prop:supplied-p126-norm-coefficients} shows that the branch trace
of $g_J$ is the pullback of the downstairs coefficient-ideal gauge.
Thus every such gauge is admissible. The density proof of
\Cref{lem:supplied-p126-gauge-density} uses only compactness and ideal
arithmetic, so it applies on $Y$ here. Closedness of $\mathcal D_f$
therefore makes every continuous function admissible. This is the
finite locally free norm of \cite[Tag~0BD2]{StacksProject}; flat going
down is \cite[Tag~00HS]{StacksProject}.
\end{proof}

For a reduced excellent curve, the conductor quotient is zero-dimensional
and its primary link is empty. Thus the normalization quotient is the
finite set of DVR branches. If $A$ is an excellent DVR with uniformizer
$\pi$ and $D$ its normalization in a finite field extension, write
$e_i=\operatorname{ord}_i(\pi)$ and
$f_i=[D/\mathfrak n_i:A/(\pi)]$. Finite freeness and the length of
$D/\pi D$ give $[\operatorname{Frac}(D):\operatorname{Frac}(A)]=\sum_ie_if_i$.
Smith normal form for multiplication gives
$\operatorname{ord}_A N(b)=\sum_i f_i\operatorname{ord}_i(b)$.
On the $\pi^r$-section the branch point is
$\operatorname{ord}_i/(re_i)$, so its trace weight is $e_if_i$.
These formulas use finite freeness over a DVR, not a defectless
hypothesis for arbitrary valuation rings.

Degree-normalized averages need not compose when intermediate ranks
vary by component. Let $A=k[[t]]$, $B=A\times A$ with diagonal map,
and $D=A\times k[[u]]$, with the second map $t\mapsto u^2$.
The direct average has weights $(1/3,2/3)$, whereas iterating the
componentwise averages gives $(1/2,1/2)$. Their values at $(0,1)$ are
$2/3$ and $1/2$. The unnormalized traces both have weights $(1,2)$
and compose, including in characteristic two.

\subsection{Changing the primary section}

For a homogeneous function $H$ of degree one on the centered cone,
let $h_{\mathfrak a}$ be its representative on $V_{\mathfrak a}$.
The exact change to the $\mathfrak b$-section is
\[
 h_{\mathfrak b}\bigl(v/v(\mathfrak b)\bigr)
                  =h_{\mathfrak a}(v)/v(\mathfrak b).
 \tag{70.7j}\label{eq:supplied-p126-radial-transfer}
\]
Every extension of $v$ has the same $\mathfrak b$-gauge. Hence the
formula commutes with trace, fiber extrema and conductor descent.
For a separately chosen source boundary $\mathfrak d$, first express
its representative on the $\mathfrak aB$-section as
$w(\mathfrak d)h_{\mathfrak d}(w/w(\mathfrak d))$.
For nonprimary gauges this conversion is confined to the domain where
the gauge is positive and finite, as in
\eqref{eq:p121-nonprimary-gauge-domain}.

The section establishes finite-map calculus on continuous valuation-link
functions, its exact conductor domain and a counterexample to a universal
continuous degree trace. Constructing independently geometric singular
psh cones and their finite-map, envelope and Monge--Amp\`ere compatibility
remains a separate question.
